\section{Introduction}

In an extended nuclear-recoil window with an exposure of $2.84$ tonne\,yr, the LUX-ZEPLIN (LZ) Collaboration reported a candidate event consistent with a recoil energy of $E_R=248\pm23_{\rm stat}\pm23_{\rm syst}$~keV. Across the models considered by LZ, the excess over the background-only hypothesis reached a local significance of $3.4\sigma$ and a global significance of $2.6\sigma$~\cite{LZ}. A dark sector that releases energy in the scattering, rather than absorbing it, can convert a moderately heavy dark state into an energetic xenon recoil; this exothermic mechanism has been studied as an explanation of both the xenon excess and the Galactic Centre gamma-ray excess~\cite{GHK}, and the specific case of a pseudo-Dirac pair coupled through a leptophobic vector mediator has been confronted with the LZ event and with collider data~\cite{DHX}.

The same interactions that produce the recoil also govern the annihilation rate in the early universe and the production of invisible final states at the LHC. A vector mediator with flavour-dependent couplings ties these three observables together: the recoil rate in xenon is controlled by the light-quark couplings coherently summed over the nucleus, the relic abundance by the annihilation channels that are kinematically open at freeze-out, and the collider rate by top-associated production, which probes the top coupling directly. When the couplings grow with quark mass, as they do in the electroweak-broken effective description adopted here, the three probes are not redundant. Xenon scattering samples the up and down content of the nucleon, freeze-out samples all open quark channels, and top-associated production samples the top coupling directly.

Motivated by the LZ recoil candidate, we ask whether this scenario can accommodate a reference xenon recoil yield while satisfying the Planck abundance and the CMS constraints. The recoil side enters as a normalisation condition rather than a binned comparison, so the question is which points of the $(M_V,g_t)$ plane reproduce a fixed recoil yield and the measured abundance, and which of those the LHC excludes at 95\% CL.

Section~\ref{sec:model} introduces the effective model. Section~\ref{sec:lz} describes the xenon normalisation, the relic-density solution and the construction of the 470-point grid. Section~\ref{sec:cms} presents the collider analysis, Section~\ref{sec:results} the exclusion region, and Section~\ref{sec:window} the stability of its boundary under the full Planck window. Sections~\ref{sec:sys} and~\ref{sec:concl} collect the uncertainties and the conclusions.
